\exercisegroup

\begin{enumerate}
\def\labelenumi{\arabic{enumi})}
\item
  设 E 与 F 是两个 Hausdorff 局部凸空间，\(E'\) 与 \(F'\) 分别是它们的对偶。证明：若对 F 的每个向量子空间 N，\(\tau(\mathrm{F}, \mathrm{F}')\) 在 N 上诱导的拓扑都与 \(\tau(\mathrm{N}, \mathrm{F}'/\mathrm{N}^{\circ})\) 相同 (IV, 第 51 页, exerc. 20)，则每个对于拓扑 \(\sigma(\mathrm{E}, \mathrm{E}')\) 与 \(\sigma(\mathrm{F}, \mathrm{F}')\) 的从 E 到 F 中的严格态射，也是对于拓扑 \(\tau(\mathrm{E}, \mathrm{E}')\) 与 \(\tau(\mathrm{F}, \mathrm{F}')\) 的严格态射（参见 IV, 第 10 页, 命题 11）。考察 F 可度量的情形。
\item
  设 E 是一个 Hausdorff 局部凸空间，使得在对偶 \(E'\) 中存在一个无限维凸集 \(B'\) ，它对于 \(\sigma(E', E)\) 是紧的（例如当 E 是一个带上 \(\sigma(E, E^{*})\) 的无限维向量空间时，该条件即成立）。证明：存在一个线性型 \(u \in E'^{*}\) ，它在 \(B'\) 上是无界的；由此推出 \(B'\) 对于拓扑 \(\sigma(E', F)\) 不是紧的，其中 F 是子空间 \(E + Ru\) （\(E'^{*}\) 的）。由此推断：从 E 到 F 的典范单射对于拓扑 \(\sigma(E, E')\) 与 \(\sigma(F, E')\) 是一个严格态射，但对于拓扑 \(\tau(E, E')\) 与 \(\tau(F, E')\) 不是。
\item
  给出一个严格单射态射 \(u\) （从 Fréchet 空间 \(E\) 到 Fréchet 空间 \(F\) 中）的例子，使得 \(^t u\) 不是从 \(F_b'\) 到 \(E_b'\) 中的严格态射（参见 IV, 第 58 页, exerc. 5, c)）。
\item
  设 E 与 F 是两个 Hausdorff 局部凸空间，u 是从 E 到 F 中的一个连续线性映射，M 是 E 的一个处处稠密的向量子空间。证明：若 u 在 M 上的限制是从 M 到 F 中的一个严格态射，则 u 是从 E 到 F 中的一个严格态射（利用 IV, 第 27 页的命题 2）。若此外还有 \(u(\mathbf{M}) = \mathbf{F}\) ，证明：对 E 中 0 的每个开的、凸且均衡的邻域 V，\(u(\mathbf{V})\) 都属于 \(\overline{u(\mathbf{V} \cap \mathbf{M})}\) 的内部。
\item
  设 \(\mathbf{E}\) 与 \(\mathbf{F}\) 是两个赋范空间，\(u\) 是从 \(\mathbf{E}\) 到 \(\mathbf{F}\) 中的一个连续线性映射。
\end{enumerate}

\begin{enumerate}
\def\labelenumi{\alph{enumi})}
\item
  证明：若 \(u\) 是从 \(E\) 到 \(F\) 中的一个严格态射，则 \({}^t u\) 是从强对偶 \(F_b'\) 到强对偶 \(E_b'\) 中的一个严格态射。
\item
  假设 E 是完备的；证明：若 \(^{t}u\) 是从 \(F_{b}^{\prime}\) 到 \(E_{b}^{\prime}\) 中的一个严格态射，则 u 是从 E 到 F 中的一个严格态射，且 \(^{t}u\) 是从 \(F^{\prime}\) 到 \(E^{\prime}\) 中的一个严格态射（对于弱拓扑 \(\sigma(F^{\prime}, F)\) 与 \(\sigma(E^{\prime}, E)\) ）（把 F 视为它的完备化的一个子空间）。
\item
  给出一个例子：E 不是完备的，F 是完备的，\(^{t}u\) 对于强拓扑以及对于弱拓扑都是从 \(F'\) 到 \(E'\) 中的一个严格单射态射，但 u 不是从 E 到 F 中的严格态射（参见 II, 第 74 页, exerc. 5）。
\item
  给出一个例子：E 不是完备的，F 是完备的，u 是从 E 到 F 中的一个严格单射态射，\(^{t}u\) 对于强拓扑是从 F' 到 E' 中的一个严格态射，但对于弱拓扑不是（取 E 在 F 中处处稠密）。
\item
  若 F 是完备的，且 \({}^t u\) 对于弱拓扑是从 \(\mathrm{F}'\) 到 \(\mathrm{E}'\) 中的一个严格态射，则 \({}^t u\) 对于强拓扑是从 \(\mathrm{F}'\) 到 \(\mathrm{E}'\) 中的一个严格态射（把 \(u\) 延拓到 \(\hat{\mathrm{E}}\) 上）。
\item
  给出一个例子：E 是完备的，F 不是完备的，\(^{t}u\) 对于弱拓扑是从 \(F'\) 到 \(E'\) 中的一个严格态射，但对于强拓扑不是（参见 II, 第 74 页, exerc. 5）。
\end{enumerate}

\begin{enumerate}
\def\labelenumi{\arabic{enumi})}
\setcounter{enumi}{5}
\item
  设 E 是可度量的局部凸空间 \(\ell^{1}(\mathbf{N})\) (I, 第 4 页) ，带上由乘积空间 \(R^{N}\) 的拓扑诱导的拓扑；它的对偶 \(E'\) 可以与 \(\mathbf{R}^{(N)}\) 等同，并且拓扑 \(\tau(\mathrm{E}',\mathrm{E})\) 是在 \(E'\) 上由 \(c_{0}(\mathbf{N})\) 的范数拓扑诱导的拓扑 (IV, 第 47 页, exerc. 1)。证明：若 u 是从 E 到一个 Hausdorff 桶型空间 F 上的满射连续线性映射，则 F 必是有限维的。（注意：\({}^{t}u\) 是从 \(F'\) （带上 \(\sigma(\mathrm{F}',\mathrm{F})\) ）到 \(E'\) （带上 \(\sigma(\mathrm{E}',\mathrm{E})\) ）的一个子空间上的同构；由 F 是桶型的这一事实，推出 \({}^{t}u(F')\) 带上由 \(\tau(\mathrm{E}',\mathrm{E})\) 诱导的拓扑是一个 Banach 空间，然后利用 II, 第 80 页的 exerc. 24。）
\item
  \begin{enumerate}
  \def\labelenumii{\alph{enumii})}
  \tightlist
  \item
    设 \(\mathbf{E}\) 是一个 Banach 空间，\((x_{\alpha})_{\alpha \in \mathbf{A}}\) 是 \(\mathbf{E}\) 的单位球面中的一个处处稠密集。设 \(u\) 是从空间 \(\ell^1 (\mathbf{A})(\mathrm{I},\mathrm{p.~4})\) 到 \(\mathbf{E}\) 中的线性映射，由 \(u(t) = \sum_{\alpha \in \mathbf{A}}t(\alpha)x_{\alpha}\) 定义。
  \end{enumerate}
\end{enumerate}

对所有 \(t = (t(\alpha))_{\alpha \in \mathbf{A}}\) （属于 \(\ell^1 (\mathbf{A})\) ）成立。证明：\(u\) 是从 \(\ell^1 (\mathbf{A})\) 到 \(\mathbf{E}\) 上的一个严格态射，从而 \(\mathbf{E}\) 同构于 \(\ell^1 (\mathbf{A})\) 的一个商空间。

\begin{enumerate}
\def\labelenumi{\alph{enumi})}
\setcounter{enumi}{1}
\tightlist
\item
  由 a) 出发，给出 \(\ell^{1}(\mathbf{N})\) 的一个在 \(\ell^{1}(\mathbf{N})\) 中没有拓扑补的闭子空间的例子（对 E 取 \(c_{0}(\mathbf{N})\) ，并利用 IV, 第 54 页, exerc. 15, b) 与第 55 页, exerc. 18）。
\end{enumerate}

\begin{enumerate}
\def\labelenumi{\arabic{enumi})}
\setcounter{enumi}{7}
\tightlist
\item
  对每个整数 \(n > 0\) ，设 \(a^{(n)}\) 是二重序列 \(a^{(n)} = (a_{ij}^{(n)})_{i\geqslant 1,j\geqslant 1}\) ，其定义为：\(a_{ij}^{(n)} = j^n\) 对每一对 \((i,j)\) （满足 \(i < n\) ），而 \(a_{ij}^{(n)} = i^n\) 对每一对 \((i,j)\) （满足 \(i\geqslant n\) ）。设 \(E\) 是由所有满足如下条件的实数二重序列 \(x = (x_{ij})\) 组成的向量空间：对每个整数 \(n > 0\) ，有 \(p_n(x) = \sum_{i,j}a_{ij}^{(n)}|x_{ij}| < +\infty\) ；诸半范数 \(p_n\) 在 E 上定义
\end{enumerate}

一个 Fréchet 空间兼 Montel 空间的拓扑 (IV, 第 60 页, exerc. 11)；E 的对偶 E' 是一个 (DF) 空间兼 Montel 空间（从而是超有界型且自反的），它等同于由所有满足如下条件的序列 \(x' = (x'_{ij})\) 组成的空间：至少对一个指标 n，有关系 \(\sup |a_{ij}^{(n)}|^{-1} |x'_{ij}| < +\infty\) (IV, 第 47 页, exerc. 1, c))。

\begin{enumerate}
\def\labelenumi{\alph{enumi})}
\item
  对每个 \(x = (x_{ij}) \in E\) ，令 \(y_j = \sum_i x_{ij}\) （对所有 \(j \geqslant 1\) ）。证明：\(\sum_j |y_j| < +\infty\) ；设 \(u(x)\) 表示序列 \((y_i) \in \ell^1(\mathbf{N})\) ；证明 u 是从 E 到 \(F = \ell^1(\mathbf{N})\) 中的一个连续线性映射，并且对每个序列 \(y' = (y'_j) \in F' = \ell^\infty(\mathbf{N})\) ，\(^t u(y')\) 是序列 \((z'_{ij}) \in E'\) ，它由 \(z'_{ij} = y'_j\) （对每个指标 i）确定。由此推出：\(^t u\) 是从 \(\ell^\infty(\mathbf{N})\) 到 \(E'\) 的一个对于 \(\sigma(E', E)\) 闭的子空间上的单射线性映射，从而 u 对于初始拓扑是从 E 到 \(\ell^1(\mathbf{N})\) 上的一个严格态射，且 \(^t u\) 是从 \(F' = \ell^\infty(\mathbf{N})\) 到 \(^t u(F')\) 上的一个同构（对于弱拓扑 \(\sigma(F', F)\) 与 \(\sigma(E', E)\) ）。
\item
  设 \(M = u^{-1}(0)\) ；则 \(M^{\circ} = {}^{t}u(F')\) 。证明：在 \({}^{t}u\) 下，\(M^{\circ}\) 上由 \(E'\) 的强拓扑诱导的拓扑的原像，是 F 的紧子集上一致收敛的拓扑 (IV, 第 28 页, 定理 1 及第 54 页, exerc. 15)。由此推出：在 \(M^{\circ}\) 上，由强拓扑 \(\beta(E', E)\) 诱导的拓扑不是强拓扑 \(\beta(M^{\circ}, E/M)\) ，并且对于由 \(\beta(E', E)\) 诱导的拓扑，\(M^{\circ}\) 不是次桶型空间，尽管它是一个超有界型 Montel 空间的闭子空间；另一方面，E/M（它是 Fréchet 空间兼 Montel 空间关于一个闭子空间的商）不是自反的。证明：E/M 中存在一些有界集，它们不是 E 的有界子集的典范像。
\end{enumerate}

\begin{enumerate}
\def\labelenumi{\arabic{enumi})}
\setcounter{enumi}{8}
\tightlist
\item
  设 A 是一个可数集。考虑三个向量空间对 \((P, P')\) ，\((Q, Q')\) ，\((E, E')\) ，这六个空间都是 \(R^{A}\) 的向量子空间，且都包含直和子空间 \(\mathbf{R}^{(A)}\) ；此外，假设对每一点 \(x \in P\) （相应地，\(x \in Q\) ，\(x \in E\) ）以及每一点 \(x' \in P'\) （相应地，\(x' \in Q'\) ，\(x' \in E'\) ），族 \((x_{\alpha}x_{\alpha}')_{\alpha \in A}\) 是可和的，并令 \(\langle x, x' \rangle = \sum_{\alpha \in A} x_{\alpha}x_{\alpha}'\) ；
\end{enumerate}

这个双线性型使 P 与 P' （相应地，Q 与 Q' ，E 与 E' ）处于分离对偶之中。

\begin{enumerate}
\def\labelenumi{\alph{enumi})}
\item
  假设 \(\mathrm{E} = \mathrm{P} \cap \mathrm{Q}\) ，\(\mathrm{E}' \supset \mathrm{P}' + \mathrm{Q}'\) 且 \(\mathrm{E}' \neq \mathrm{P}' + \mathrm{Q}'\) 。考虑线性映射 \(u: x \mapsto (x, x)\) （从 \(\mathrm{E}\) 到 \(\mathrm{F} = \mathrm{P} \times \mathrm{Q}\) 中），它与 \(\mathrm{F}' = \mathrm{P}' \times \mathrm{Q}'\) 处于分离对偶之中。证明：\(u\) 对于弱拓扑 \(\sigma(\mathrm{E}, \mathrm{E}')\) 与 \(\sigma(\mathrm{F}, \mathrm{F}')\) 是连续的，并且它的像 \(\mathbf{M} = u(\mathbf{E})\) 对于 \(\sigma(\mathbf{F}, \mathbf{F}')\) 是一个闭子空间；由此推出：\(^t u\) 是从 \(\mathbf{F}'\) 到 \(\mathbf{E}'\) 中的一个严格态射（对于拓扑 \(\sigma(\mathbf{F}', \mathbf{F})\) 与 \(\sigma(\mathbf{E}', \mathbf{E})\) ），而 \(\mathbf{N} = {}^{t}u(\mathbf{F}')\) 在 \(\mathbf{E}'\) 中对于 \(\sigma(\mathbf{E}', \mathbf{E})\) 不是闭的。若 \(\mathbf{E}'\) 对于拓扑 \(\tau(\mathbf{E}', \mathbf{E})\) 是可度量的，推出：\(^t u\) 也是从 \(\mathbf{F}'\) 到 \(\mathbf{E}'\) 中的一个严格态射（对于拓扑 \(\tau(\mathbf{F}', \mathbf{F})\) 与 \(\tau(\mathbf{E}', \mathbf{E})\) ）。
\item
  此外，假设 \(E'\) 对于 \(\tau(E', E)\) 是一个 Fréchet 空间。则 \(F'/M^{\circ}\) 带上 \(\tau(F', F)\) 关于 \(M^{\circ}\) 的商拓扑后不是半完备的，并且 \(F'/M^{\circ}\) 中存在一些对于 \(\tau(F'/M^{\circ}, M)\) 不是相对紧的有界集。
\item
  在与 b) 相同的假设下，设 \(x'\) 是 \(E'\) 的一个不属于 \(\mathbf{N} = {}^{u}(\mathbf{F}')\) 的元素；对每个 \(y \in M\) ，令 \(v(y) = \langle x, x' \rangle\) ，其中 \(x \in E\) 是使 \(u(x) = y\) 的唯一元素。证明：M 上的线性型 v 对于拓扑 \(\sigma(F, F')\) 不是连续的，但它在 M 的每个有界子集上的限制对于 \(\sigma(M, F'/M^{\circ})\) 是连续的。由此推出：\(L = v^{-1}(0)\) 是 F 的一个向量子空间，它与 F 的每个有界闭子集（对于 \(\sigma(F, F')\) ）的交对于 \(\sigma(F, F')\) 是闭的，但它本身在 F 中对于 \(\sigma(F, F')\) 不是闭的。
\end{enumerate}

¶ 10) a) 设 G、H 是两个自反 Banach 空间，满足 \(\mathbf{R}^{(\mathbf{N})} \subset \mathbf{G} \subset \mathbf{H} \subset \mathbf{R}^{\mathbf{N}}\) （* 例如 \(\mathbf{G} = \ell^{r}(\mathbf{N})\) 且 \(\mathbf{H} = \ell^{p}(\mathbf{N})\) ，其中 \(1 < r < p < +\infty\) ）。取 \(\mathbf{A} = \mathbf{N} \times \mathbf{N}\) ，\(\mathbf{P} = \mathbf{H}^{(\mathbf{N})}\) （拓扑直和），\(\mathbf{P}' = \mathbf{H}'^{\mathbf{N}}\) ，\(\mathbf{Q} = \mathbf{G}^{\mathbf{N}}\) ，\(\mathbf{Q}' = \mathbf{G}'^{(\mathbf{N})}\) ，\(\mathbf{E} = \mathbf{G}^{(\mathbf{N})}\) ，\(\mathbf{E}' = \mathbf{G}'^{\mathbf{N}}\) ，证明：exerc. 9, a) 的条件满足；\(\mathbf{E}'\) 是一个自反 Fréchet 空间，且 E、F、F' 都是自反 Fréchet 空间的严格归纳极限（从而是完备且自反的）。

\begin{enumerate}
\def\labelenumi{\alph{enumi})}
\setcounter{enumi}{1}
\tightlist
\item
  由 a) 与 exerc. 9 出发，构造具有下列性质的例子：
\end{enumerate}

\(\alpha)\) 自反 Fréchet 空间的严格归纳极限的一个商空间，它既非拟完备亦非半自反。

β) 自反 Fréchet 空间的严格归纳极限的一个闭子空间，它不是自反的，且其对偶不是完备的。

\(\gamma)\) 自反 Fréchet 空间的严格归纳极限 \(\mathbf{E}_n\) 的一个向量子空间，它不是闭的（从而不是桶型的），但它与每个子空间 \(\mathbf{E}_n\) 的交都是闭的。

δ) 自反 Fréchet 空间的严格归纳极限的对偶的一个非闭子空间，它与每个弱紧子集的交都是弱紧的（参见 IV, 第 25 页, 推论 2）。

\[
\mathrm{E} _ {n}
\]

\[
\mathrm{F} _ {n} (c f.
\]

\[
\mathrm{F} _ {n} = \overline {{\mathrm{F} \cap \mathrm{E} _ {n}}}
\]

\[
u _ {n}
\]

\[
v _ {n}
\]

\[
\mathrm{F} \cap \mathrm{E} _ {n}
\]

\[
\overline {{\mathrm{F} \cap \mathrm{E} _ {n}}}
\]

\[
v _ {n}
\]

\begin{enumerate}
\def\labelenumi{\alph{enumi})}
\setcounter{enumi}{1}
\tightlist
\item
  假设 F 在 E 中不是闭的，但对所有 n，\(F \cap E_{n}\) 在 \(E_{n}\) 中是闭的 (exerc. 10, b))。设 \(A_{n}\) 是 \(E_{n}\) 中的一个可数稠密集，G 是 E 中由 F 与诸 \(A_{n}\) 的并所生成的向量子空间。证明：G 是一个有界型空间，其中 F 是一个具有可数基的补的子空间，但 F 不是次桶型的（参见 III, 第 41 页, § 2, exerc. 4 及第 45 页, exerc. 17）。
\end{enumerate}

\begin{enumerate}
\def\labelenumi{\arabic{enumi})}
\setcounter{enumi}{11}
\tightlist
\item
  设 E、F 是两个 Fréchet 空间，u 是从 E 到 F 中的一个严格态射。证明：对每个从 E 到 F 中的有限秩连续线性映射 v，\(u + v\) 都是从 E 到 F 中的一个严格态射。
\end{enumerate}

¶ 13) a) 设 E 是一个 Hausdorff 且完备的局部凸空间。假设存在 E 中闭向量子空间的一个递减序列 ( \(F_{n}\) )，使得对 E 中 0 的每个邻域 V，存在 n 使得 \(F_{n} \subset V\) 。证明：空间 E 是极小型的 (II, 第 85 页, exerc. 13)。

\begin{enumerate}
\def\labelenumi{\alph{enumi})}
\setcounter{enumi}{1}
\item
  设 E 是一个满足第一可数公理且非极小型的 Fréchet 空间。证明：E 中存在两个闭向量子空间 M、N，使得 \(M \cap N = \{0\}\) 且 \(M + N\) 不是闭的。（设 \((x_{n}')\) 是 E 上线性无关的连续线性型的一个序列，构成 \(\sigma(E', E)\) 下的一个完全集 (III, 第 19 页, 推论 2)；设 \(L_{n}\) 是 E 中与诸 \(x_{i}'\) （指标 \(i \leqslant 2n\) ）正交的子空间；设 \(x_{n}, y_{n}\) 是 \(L_{n+1}\) 关于 \(L_{n}\) 的余中的两个线性无关向量。设 d 是定义 E 的拓扑的一个平移不变距离。利用 a)，证明存在一个数 \(\alpha > 0\) ，使得可取 \(d(0, x_{n}) \geqslant \alpha\) ，\(d(0, y_{n}) \geqslant \alpha\) 且 \(d(x_{n}, y_{n}) \leqslant 1/n\) 。证明：若 M（相应地，N）是由诸 \(x_{n}\) （相应地，\(y_{n}\) ）生成的闭向量子空间，则 M 与 N 具有所要求的性质；利用 I, 第 19 页, 推论 4。）
\item
  设 \(\mathbf{E}\) 是赋范空间的一个乘积 \(\prod_{\alpha \in \mathbf{A}} \mathrm{F}_{\alpha}\) 的闭子空间；证明：若 E 中由点的可数族生成的每个闭
\end{enumerate}

子空间都是极小型的，则 E 本身是极小型的。（用反证法，假设 E 在每个 \(F_{\alpha}\) 上的投影都等于 \(F_{\alpha}\) ，且对某个 \(\alpha\) ，\(F_{\alpha}\) 是无限维的。）

\begin{enumerate}
\def\labelenumi{\alph{enumi})}
\setcounter{enumi}{3}
\tightlist
\item
  设 E 是一个 Hausdorff 且完备的局部凸空间，其中由点的可数族生成的每个子空间都是可度量的。要使 E 是极小型的，必须且只需对 E 的每对满足 \(M \cap N = \{0\}\) 的闭向量子空间 M、N，\(M + N\) 在 E 中是闭的（利用 b) 与 c)）。
\end{enumerate}

\begin{enumerate}
\def\labelenumi{\arabic{enumi})}
\setcounter{enumi}{13}
\tightlist
\item
  \begin{enumerate}
  \def\labelenumii{\alph{enumii})}
  \tightlist
  \item
    设 E 是一个不存在连续范数的 Fréchet 空间。证明：E 中存在一个极小型的闭向量子空间 (II, 第 85 页, exerc. 13)，它是无限维的，从而在 E 中具有一个拓扑补。（存在 0 的凸、均衡邻域的一个严格递减基本序列 \((\mathrm{V}_{n})\) ，以及 E 的点的一个序列 \((x_{n})\) ，使得 \(x_{n} \notin V_{n+1}\) ，且过 \(x_{n}\) 的直线包含于 \(V_{n}\) 中，诸 \(x_{n}\) 线性无关；证明所求的空间是由诸 \(x_{n}\) 生成的闭向量子空间。）
  \end{enumerate}
\end{enumerate}

\begin{enumerate}
\def\labelenumi{\alph{enumi})}
\setcounter{enumi}{1}
\item
  设 E 是一个 Fréchet 空间，其拓扑不能由单独一个范数、而由范数的一个递增序列 \((p_{n})\) 定义。设 \(V_{n}\) 是由 \(p_{n}(x) \leqslant 1\) 定义的邻域，并设 \(A_{n}^{\prime} = V_{n}^{\circ}\) 是它在 \(E^{\prime}\) 中的极集；可以假设 \(A_{n+1}^{\prime}\) 不包含于 \(E^{\prime}\) 的、由 \(A_{n}^{\prime}\) 生成的向量子空间中 (IV, 第 49 页, exerc. 12, c))。设 \((x_{n}^{\prime})\) 是 \(E^{\prime}\) 的点的一个序列，满足 \(x_{n}^{\prime} \in A_{n}^{\prime}\) 且 \(x_{n}^{\prime}\) 不属于由 \(A_{n-1}^{\prime}\) 生成的向量子空间；证明：向量子空间 \(M^{\prime}\) （由序列 \((x_{n}^{\prime})\) 生成）在 \(E^{\prime}\) 中是弱闭的（参见 IV, 第 25 页, 推论 2），并且在 \(E^{\prime}\) 中对于拓扑 \(\sigma(E^{\prime}, E)\) 没有拓扑补。（注意：若存在一个弱连续投影 \(u^{\prime}\) （从 \(E^{\prime}\) 到 \(M^{\prime}\) 上），则由 Baire 定理，\(u^{\prime}(A_{1}^{\prime})\) 将包含于诸集 \(M^{\prime} \cap A_{n}^{\prime}\) 之一中，由此导出矛盾，因为 \(A_{1}^{\prime}\) 在 \(E^{\prime}\) 中是弱完全的。）由此推出：E 的子空间 \(M^{\circ}\) 在 E 中没有拓扑补。
\item
  设 E 是一个其拓扑不能由单独一个范数定义的 Fréchet 空间。证明：若 E 不同构于 Banach 空间的一个乘积，则 E 中存在一个没有拓扑补的闭向量子空间。（用反证法；设 \((p_{n})_{n\geqslant1}\) 是 E 上定义 E 的拓扑的半范数的一个递增序列；设 \(F_{n}=p_{n}^{-1}(0)\) ，并设 \(E_{n+1}\) 是 \(F_{n+1}\) 关于 \(F_{n}\) 的一个拓扑补（令 \(E=F_{0}\) ）；利用 b)，证明 \(E_{n+1}\) 是一个 Banach 空间，且 E 同构于诸 \(E_{n}(n\geqslant1)\) 的乘积；为此利用 I, 第 17 页, 定理 1。）
\end{enumerate}

\begin{enumerate}
\def\labelenumi{\arabic{enumi})}
\setcounter{enumi}{14}
\tightlist
\item
  \begin{enumerate}
  \def\labelenumii{\alph{enumii})}
  \tightlist
  \item
    设 E 是一个无限维赋范空间（实的或复的）。证明：E 中存在一个序列 \((x_{n})\) ，使得对标量的每个有界序列 \((\lambda_{n})\) ，存在 E 上的一个连续线性型 \(x^{\prime}\) ，使得 \(\langle x_n,x'\rangle = \lambda_n\) 对所有 \(n\) 成立。（构造 E 的对偶 E' 中点的一个序列 \((x_n')\) 以及 E 中点的一个序列 \((x_{n})\) ，使得 \(\langle x_i,x_j'\rangle = \delta_{ij}\) 且 \(\| x_n'\| \leqslant 2^{-n}\) 。）
  \end{enumerate}
\end{enumerate}

\begin{enumerate}
\def\labelenumi{\alph{enumi})}
\setcounter{enumi}{1}
\tightlist
\item
  要使可度量的局部凸空间 \(\mathbf{E}\) 具有 \(a\) ) 中所述的性质，必须且只需 \(\mathbf{E}\) 的完备化不是极小型空间 (II, 第 85 页, exerc. 13)。
\end{enumerate}

\begin{enumerate}
\def\labelenumi{\arabic{enumi})}
\setcounter{enumi}{15}
\tightlist
\item
  \begin{enumerate}
  \def\labelenumii{\alph{enumii})}
  \tightlist
  \item
    设 E 与 F 是两个无限维复赋范空间，u 是从 E 到 F 上的一个双射半线性映射，它相对于 C 的一个自同构 \(\sigma\) ，并且把 E 的每个闭超平面变成 F 的闭超平面。证明：C 的自同构 \(\sigma\) 必是连续的（从而或者是恒同，或者是自同构 \(\xi \mapsto \overline{\xi}\) ）。（用反证法：设 \((x_{n})\) 是 E 中满足 exerc. 15, a) 的条件的点的一个序列，并设 \((\lambda_{n})\) 是复数的一个有界序列，使得对所有 n 有 \(|\lambda_{n}^{\sigma}| \geqslant n\) . \(\|u(x_{n})\|\) ；若 \(x' \in E'\) 使得对所有 n 有 \(\langle x_{n}, x' \rangle = \lambda_{n}\) ，考虑 u 在闭超平面（由所有 \(x \in E\) （满足 \(\langle x, x' \rangle = 1\) ）组成）上的像，并导出矛盾。）
  \end{enumerate}
\end{enumerate}

\begin{enumerate}
\def\labelenumi{\alph{enumi})}
\setcounter{enumi}{1}
\tightlist
\item
  由 a) 推出：u 是从 E 到 F 中的一个连续半线性映射 (IV, 第 7 页, 推论)。
\item
  设 \(\sigma\) 是域 C 的一个不连续自同构。证明：双射 \((\xi_{n})\mapsto(\xi_{n}^{\sigma})\) （从 \(C^{N}\) 到其自身上的）把每个闭超平面变成闭超平面（参见 IV, 第 14 页, 命题 15）。
\end{enumerate}

17) 设 E、F 是两个 Banach 空间，u 是从 E 到 F 上的一个严格态射。则存在
一个数 \(m > 0\) ，使得对每个 \(\varepsilon \in ]0, m[\) 以及每一点 \(y \in F\) ，存在 \(z \in E\) ，使
得 \(u(z) = y\) 且 \(\| z \| \leqslant (m - \varepsilon)^{-1} \| y \|\)。
a) 设 \(\mathbf{B}\) 是 E 中的一个闭球 \(\| x - a \| \leqslant r\) ，并设 w 是从 \(\mathbf{B}\) 到 F 中的一个映射，满足
下列条件：\(1^{\circ}\ \sup_{x \in \mathbf{B}} \| w(x) \| = M < +\infty; 2^{\circ}\) 存在一个数 \(k > 0\) ，
使得对所有 \(x, x'\) （属于 \(\mathbf{B}\) ），有 \(\| w(x) - w(x') \| \leqslant k \| x - x' \|\) （“Lipschitz 条件”）。证明：若 \(k < m\) 且 \(M < r(m - k)\) ，则 \(\mathbf{B}\) 在映射 \(x \mapsto u(x) + w(x)\)
下的像包含一个以 \(b = u(a)\) 为中心的球。（证明：对每个 \(y \in F\) （足够接近于 \(b\) 的），可以定义
一个序列 \((x_n)_{n \geqslant 0}\) （\(\mathbf{B}\) 的点的），使得 \(u(x_0) = y\) 且 \(u(x_n) = y - w(x_{n-1})\) （对 \(n \geqslant 1\) ），
并且使得 \((x_n)\) 收敛于 \(\mathbf{B}\) 的一点。）


\begin{enumerate}
\def\labelenumi{\alph{enumi})}
\setcounter{enumi}{1}
\tightlist
\item
  假设 \(w\) 是从整个 \(E\) 到 \(F\) 中的一个映射，使得 \(\| w(x) - w(x') \| \leqslant k \| x - x' \|\) 对所有 \(x, x'\) （属于 \(E\) ）成立。证明：若 \(k < m\) ，则 \(u + w\) 是从 \(E\) 到 \(F\) 上的一个满射且开的映射。c) 特别地，若 \(w\) 是从 \(E\) 到 \(F\) 中的一个连续线性映射且 \(\| w \| < m\) ，则 \(v = u + w\) 是从 \(E\) 到 \(F\) 上的一个严格态射。若 \(\varepsilon \in ]0, m[\) ，则对每个 \(x \in v^{-1}(0)\) （满足 \(\| x \| = 1\) ），存在 \(x_0 \in u^{-1}(0)\) 使得 \(\| x - x_0 \| \leqslant \| w \| / (m - \varepsilon)\) ；若此外 \(\| w \| < m - \varepsilon\) ，则对每个 \(x_0 \in u^{-1}(0)\) （满足 \(\| x_0 \| = 1\) ），存在 \(x \in v^{-1}(0)\) 使得
\end{enumerate}

\[
\left\| x - x _ {0} \right\| \leqslant \| w \| / (m - \varepsilon - \| w \|).
\]

由此推出：若存在 E 的一个闭向量子空间 G，它是 \(u^{-1}(0)\) 的补，则 G 也是 \(v^{-1}(0)\) 的补，只要 \(\|w\|\) 足够小（用反证法：证明 \(v^{-1}(0)\) 到 \(u^{-1}(0)\) 上的投影不能包含于 \(u^{-1}(0)\) 的一个闭超平面中；另一方面，注意存在 a > 0，使得每个点 \(x \in G\) （满足 \(\|x\| = 1\) ）都有 \(\geqslant a\) 的距离到 \(u^{-1}(0)\) ）。

\begin{enumerate}
\def\labelenumi{\arabic{enumi})}
\setcounter{enumi}{17}
\tightlist
\item
  \begin{enumerate}
  \def\labelenumii{\alph{enumii})}
  \tightlist
  \item
    设 E、F 是两个赋范空间，u 是从 E 到 F 中的一个严格单射态射；则存在一个数 m > 0，使得 \(\|u(x)\| \geqslant m \|x\|\) 对所有 \(x \in E\) 成立。证明：若 w 是从 E 到 F 中的一个连续线性映射且 \(\|w\| < m\) ，则 \(v = u + w\) 是从 E 到 F 中的一个严格单射态射。此外，对每个 \(y_0 \in u(E)\) （满足 \(\|y_0\| = 1\) ），存在 \(y \in v(E)\) 使得 \(\|y - y_0\| \leqslant \|w\| / m\) ；对每个 \(y \in v(E)\) （满足 \(\|y\| = 1\) ），存在 \(y_0 \in u(E)\) 使得 \(\|y - y_0\| \leqslant \|w\| / (m - \|w\|)\) 。
  \end{enumerate}
\end{enumerate}

\begin{enumerate}
\def\labelenumi{\alph{enumi})}
\setcounter{enumi}{1}
\tightlist
\item
  由 a) 推出：若此外 E 与 F 都是 Banach 空间，且存在 F 的一个闭向量子空间 G，它是闭子空间 \(u(\mathrm{E})\) 的补，则 G 也是 \(v(\mathrm{E})\) 的补，只要 \(\|w\|\) 足够小（按 exerc. 17, c) 的方式论证）。
\end{enumerate}

\begin{enumerate}
\def\labelenumi{\arabic{enumi})}
\setcounter{enumi}{18}
\item
  设 E 是 Banach 空间 \(\mathcal{C}(\left[-1,1\right];\mathbf{R})\) （由 \([-1,1]\) 到 R 中的全体连续映射组成）中由多项式组成的子空间。类似地，设 F 是 \(\mathcal{C}(\left[0,1\right];\mathbf{R})\) 中由全体多项式组成的子空间。设 u 是使每个多项式 \(f\in E\) 对应于 F 中的多项式 \(t\mapsto\frac{1}{2}(f(\sqrt{t})+f(-\sqrt{t}))\) 的映射。又设 w 是从 E 到 F 中把每个多项式 \(f\in E\) 对应于它在 \([0,1]\) 上的限制的映射。证明：u 是从 E 到 F 上的一个严格态射，但对每个 \(\varepsilon>0\) ，\(u+\varepsilon w\) 都不是从 E 到 F 中的严格态射。
\item
  \begin{enumerate}
  \def\labelenumii{\alph{enumii})}
  \tightlist
  \item
    设 E、F 是两个 Banach 空间，u 是从 E 到 F 中的一个严格态射，且 \(u^{-1}(0)\) 是有限维的。证明：对每个范数足够小的从 E 到 F 中的连续线性映射 w，\(v = u + w\) 都是从 E 到 F 中的一个严格态射，且 \(\dim v^{-1}(0) \leqslant \dim u^{-1}(0)\) 。（把 E 写成 \(u^{-1}(0)\) 与一个闭子空间的拓扑直和，并利用 IV, 第 66 页的 exerc. 18。）
  \end{enumerate}
\end{enumerate}

\begin{enumerate}
\def\labelenumi{\alph{enumi})}
\setcounter{enumi}{1}
\tightlist
\item
  设 E、F 是两个 Banach 空间，u 是从 E 到 F 中的一个连续线性映射，使得 \(u(\mathrm{E})\) 在 F 中有有限余维数。则 \(u(\mathrm{E})\) 是闭的，且 u 是一个严格态射 (I, 第 28 页, exerc. 4)。证明：对每个范数足够小的从 E 到 F 中的连续线性映射 w，\(v = u + w\) 都是从 E 到 F 中的一个严格态射，并且
\end{enumerate}

\[
\operatorname{codim} (v (\mathrm{E})) \leqslant \operatorname{codim} (u (\mathrm{E}))
\]

（考虑 \({}^t v = {}^t u + {}^t w\) ，并利用 a) 与 IV, 第 30 页, 推论 3）。

\begin{enumerate}
\def\labelenumi{\arabic{enumi})}
\setcounter{enumi}{20}
\tightlist
\item
  设 E 与 F 是两个 Banach 空间。我们称从 E 到 F 中的一个连续线性映射为 Fredholm 算子（或拟同构），如果 \(u^{-1}(0)\) 是有限维的且 \(u(\mathrm{E})\) 有有限余维数（这蕴含 \(u(\mathrm{E})\) 在 F 中是闭的且 u 是一个严格态射）；数 \(\operatorname{Ind}(u) = \operatorname{codim}(u(\mathrm{E})) - \dim(u^{-1}(0))\) 称为 u 的指标。
\end{enumerate}

\begin{enumerate}
\def\labelenumi{\alph{enumi})}
\item
  证明：\({}^t u: \mathrm{F}' \to \mathrm{E}'\) 也是一个 Fredholm 算子，并且 \(\operatorname{Ind}({}^t u) = -\operatorname{Ind}(u)\) 。
\item
  若 \(u: \mathrm{E} \to \mathrm{F}\) 与 \(v: \mathrm{F} \to \mathrm{G}\) 是两个 Fredholm 算子，则 \(v \circ u: \mathrm{E} \to \mathrm{G}\) 也是，且 \(\operatorname{Ind}(v \circ u) = \operatorname{Ind}(u) + \operatorname{Ind}(v)\) 。
\item
  若 \(w: E \to F\) 是一个有限秩或范数足够小的连续线性映射，则 \(u + w\) 是一个 Fredholm 算子，且 \(\operatorname{Ind}(u + w) = \operatorname{Ind}(u)\) （利用 IV, 第 65 页的 exerc. 17, c) 及 18, b)）。
\end{enumerate}

\begin{enumerate}
\def\labelenumi{\arabic{enumi})}
\setcounter{enumi}{21}
\tightlist
\item
  设 \(\mathbf{X}\) 是一个实 Banach 空间，\(\mathbf{E}\) 是 \(\mathbf{X}\) 的一个有限维子空间。
\end{enumerate}

\begin{enumerate}
\def\labelenumi{\alph{enumi})}
\item
  设 S 是 X 中的单位球面。证明：对每个 \(\varepsilon > 0\) ，存在有限多个线性无关的点 \(z_{i} \in S\) （ \(1 \leqslant i \leqslant r\) ），使得对每个 \(x \in S \cap E\) ，存在一个指标 \(i\) 使得 \(\| x - z_{i} \| \leqslant \varepsilon\) 。
\item
  设 \(z_{i}^{\prime}(1 \leqslant i \leqslant r)\) 是对偶 \(E^{\prime}\) （\(E\) 的对偶）中的点，满足 \(\| z_{i}^{\prime}\| \geqslant 1\) 对每个 \(i\) ，且 \(\langle z_i, z_j' \rangle = \delta_{ij}\) （Kronecker 符号），并设 \(F\) 是余维数为 \(r\) 的闭子空间（\(E\) 中的），它与 \(\mathbf{E}'\) 中由诸 \(z_{i}^{\prime}\) 生成的子空间正交。证明：对每个 \(x\in \mathrm{S}\cap \mathrm{E}\) 以及每个 \(y\in \mathrm{F}\) ，\(\| x + y\| \geqslant 1 - \varepsilon\) ，特别地 \(\mathrm{E}\cap \mathrm{F} = \{0\}\) ，从而和 \(\mathrm{E} + \mathrm{F}\) 是拓扑直和。
\item
  由 \(b)\) 推出：连续投影 \(P\) （从 \(\mathrm{E} + \mathrm{F}\) 到 \(\mathrm{E}\) 上的、对应于拓扑直和分解 \(\mathrm{E} \oplus \mathrm{F}\) 的）的范数 \(\| P \| < 1 / (1 - \varepsilon)\) 。
\end{enumerate}

\begin{enumerate}
\def\labelenumi{\arabic{enumi})}
\setcounter{enumi}{22}
\tightlist
\item
  设 \(X\) 是一个无限维实 Banach 空间。
\end{enumerate}

\begin{enumerate}
\def\labelenumi{\alph{enumi})}
\item
  假设对每个 \(\lambda > 0\) ，存在 X 的一个有限维子空间 \(\mathrm{E}_{\lambda}\) ，使得不存在 X 上的连续投影 \(P\) ，它以 \(\mathrm{E}_{\lambda}\) 为像且满足 \(\| P \| \leqslant \lambda\) 。证明：X 的每个有限余维数的闭子空间 Y 都具有与 X 相同的性质。
\item
  对 n 用归纳法证明：存在 X 的闭子空间（均有有限余维数）的一个递减序列 \((\mathrm{X}_{n})\) ，并且对每个 n，存在 X 的一个有限维子空间 \(E_{n}\) ，使得：\(1^{\circ}\) 和 \(E_{1} + E_{2} + \cdots + E_{n}\) 是直和，且存在一个连续投影 \(P_{n}\) （\(X_{n}\) 上的，范数 \(\leqslant 2\) ），其像为 \(E_{1} + E_{2} + \cdots + E_{n}\) ；\(2^{\circ}\) 空间 \(E_{n}\) 包含于 \((\mathrm{I} - P_{n-1})(\mathrm{X}_{n-1})\) 中；\(3^{\circ}\) 不存在 X 上以 \(E_{n}\) 为像且范数 \(\leqslant n + 2\) 的连续投影。
\item
  设 \(\bar{Z}\) 是 \(X\) 中由诸 \(E_{n}\) 的并生成的闭子空间。证明：\(Z\) 在 \(X\) 中不存在拓扑补（注意：若 \(Q\) 是 \(X\) 上的一个连续投影，以 \(\bar{Z}\) 为像，则 \((I - P_{n-1})P_nQ\) 将是 \(X\) 上以 \(E_{n}\) 为像的一个连续投影）。
\end{enumerate}
% semantic-fragments: legacy-input-seam
\relax{}
